\documentclass{beamer}
\usepackage{amsmath, booktabs}
\usetheme{Madrid}


\title{Measurement of the Level of Poverty}
\subtitle{Contemporary Economic Issues: ECON 204}
\author{Muhebullah Karimzada}
\date{Winter 2025}





\begin{document}

\frame{\titlepage}

\begin{frame}{Introduction}
    \begin{itemize}
        \item Poverty measurement is essential for assessing economic deprivation and informing policy decisions.
        \item Two broad approaches:
        \begin{enumerate}
            \item \textbf{Absolute Poverty}: Defined as a fixed threshold, such as living on less than \$2.15 per day (PPP-adjusted). This approach considers the minimum income required for basic human needs.
            \item \textbf{Relative Poverty}: Defined in relation to the overall income distribution of a society. Individuals are considered poor if they earn significantly less than the median income.
        \end{enumerate}
    \end{itemize}
\end{frame}

\begin{frame}{Absolute Poverty Measures}
    \textbf{Definition:} Absolute poverty measures assess whether individuals or households fall below a fixed income threshold, regardless of societal income levels.\\
    \textbf{Example: Poverty Line}\\
    Suppose a country sets a poverty line at \$1,000 per year per person.\\
    Individuals earning below this threshold are considered poor.
\end{frame}

\begin{frame}{Headcount Ratio (HCR)}
    \textbf{Definition:} The Headcount Ratio (HCR) measures the proportion of the population living below the poverty line.\\
    \textbf{Formula:}
    \begin{equation}
        HCR = \frac{q}{N}
    \end{equation}
    where:
    \begin{itemize}
        \item \( q \) = number of people below the poverty line
        \item \( N \) = total population
    \end{itemize}
    \textbf{Example Calculation:}
    \begin{itemize}
        \item Population: 10,000
        \item Poor individuals: 2,500
        \item \( HCR = \frac{2,500}{10,000} = 0.25 \) or 25\%
    \end{itemize}
\end{frame}

\begin{frame}{Poverty Gap Index (PGI)}
    \textbf{Definition:} The Poverty Gap Index measures the intensity of poverty by calculating how far below the poverty line poor individuals are, on average.
    
    \textbf{Formula:}
    \begin{equation}
        PGI = \frac{1}{N} \sum_{i=1}^{q} \frac{(Z - Y_i)}{Z}
    \end{equation}
    where:
    \begin{itemize}
        \item \( Z \) = poverty line
        \item \( Y_i \) = income of individual \( i \)
    \end{itemize}
    \textbf{Example:}
    \begin{itemize}
        \item \( Z = 1,000 \), \( Y_A = 800 \), \( Y_B = 700 \), \( Y_C = 500 \)
        \item \( PGI = \frac{(0.20 + 0.30 + 0.50)}{3} = 0.333 \) or 33.3\%
    \end{itemize}
\end{frame}

\begin{frame}{Multidimensional Poverty Index (MPI)}
    \textbf{Definition:} The MPI assesses poverty using multiple indicators, including education, health, and living standards. A person is considered \textbf{multidimensionally poor} if they are deprived in multiple dimensions simultaneously. Each indicator represents an essential aspect of well-being, and being deprived in multiple indicators suggests a deeper level of poverty.\\
    \textbf{Example Data:}
    \begin{table}[]
        \centering
        \begin{tabular}{lccc}
        \toprule
        \textbf{Person} & \textbf{Health} & \textbf{Education} & \textbf{Living Standards} \\
        \midrule
        A & 1 & 0 & 1 \\
        B & 0 & 1 & 0 \\
        C & 1 & 1 & 1 \\
        \bottomrule
        \end{tabular}
    \end{table}

\end{frame}




\begin{frame}{Multidimensional Poverty Index (MPI)}

    \textbf{Interpretation:}
    \begin{itemize}
        \item A value of \(1\) in an indicator means the person is deprived in that dimension.
        \item A value of \(0\) means they are not deprived.
        \item The \textbf{MPI Poverty Threshold} is typically set at 0.33, meaning a person must be deprived in at least one-third of the total indicators to be classified as multidimensionally poor.
        \item In this example, \textbf{Person A and C} are multidimensionally poor since they experience deprivation in multiple dimensions.
    \end{itemize}
\end{frame}









\begin{frame}{Policy Recommendations}
    \begin{itemize}
        \item Implement targeted cash transfer programs to alleviate extreme poverty.
        \item Improve access to education and healthcare to enhance human capital.
        \item Support microfinance initiatives for self-employment opportunities.
        \item Strengthen labor market policies to ensure fair wages and employment security.
    \end{itemize}
\end{frame}

\begin{frame}{Conclusion}
    \begin{itemize}
        \item Measuring poverty requires both absolute and relative approaches.
        \item Combining income-based and multidimensional measures provides a fuller picture.
        \item Effective policies require a deep understanding of poverty severity and distribution.
    \end{itemize}
\end{frame}


\begin{frame}
\frametitle{Problem 1: Poverty Line Calculation}
\textbf{Question:}
The poverty line in Country X is set at \$1,200 per month for a household of four people. The average income of a family in this country is \$1,800 per month. Calculate the poverty gap ratio for the following families:
\begin{enumerate}
    \item Family A: Monthly income = \$1,100
    \item Family B: Monthly income = \$1,400
\end{enumerate}



\end{frame}


\begin{frame}
\frametitle{Problem 1: Poverty Line Calculation}

\textbf{Solution:}
The poverty gap ratio for a family is calculated as:
\[
\text{Poverty Gap Ratio} = \frac{\text{Poverty Line} - \text{Income of the Household}}{\text{Poverty Line}}
\]

\textbf{Family A:}  
Income = \$1,100, Poverty Line = \$1,200  
Poverty Gap = \$1,200 - \$1,100 = \$100  
Poverty Gap Ratio = \(\frac{100}{1200} = 0.0833\)

\textbf{Family B:}  
Income = \$1,400, Poverty Line = \$1,200  
Poverty Gap = \$1,200 - \$1,400 = -\$200 \quad (\text{no poverty gap as income is above poverty line})  
Poverty Gap Ratio = 0

\end{frame}







\begin{frame}
\frametitle{Problem 2: Headcount Ratio}
\textbf{Question:}
In Country Y, the poverty line is set at \$10,000 per year. A survey of 100 households reveals the following incomes (in dollars per year):
\begin{itemize}
    \item 40 households earn less than \$10,000 per year.
    \item 60 households earn more than \$10,000 per year.
\end{itemize}
Calculate the \textbf{headcount ratio} of poverty.



\end{frame}


\begin{frame}
\frametitle{Problem 2: Headcount Ratio}

\textbf{Solution:}
The headcount ratio is calculated as the proportion of people living below the poverty line:
\[
\text{Headcount Ratio} = \frac{\text{Number of Poor Households}}{\text{Total Number of Households}}
\]
Here, the number of poor households = 40, total households = 100.

\[
\text{Headcount Ratio} = \frac{40}{100} = 0.40
\]
Thus, the headcount ratio is 0.40 or 40\%.

\end{frame}

\begin{frame}
\frametitle{Problem 3: Poverty Gap Index}
\textbf{Question:}
Consider a country with the following data:
\begin{itemize}
    \item The poverty line is \$5,000 per year.
    \item There are 100 households with the following incomes (in dollars per year):
    \begin{itemize}
        \item 20 households earn \$2,000 each.
        \item 30 households earn \$3,000 each.
        \item 50 households earn \$6,000 each.
    \end{itemize}
\end{itemize}
Calculate the \textbf{poverty gap index} for this country.


\end{frame}


\begin{frame}
\frametitle{Problem 3: Poverty Gap Index}

\textbf{Solution:}
The poverty gap for each household group is calculated as the difference between the poverty line and the income of the households. The poverty gap index is given by:
\[
\text{Poverty Gap Index} = \frac{\sum (\text{Poverty Line} - \text{Income}) \times \text{Number of Poor}}{\text{Total Income of All Households}}
\]

For the 20 households earning \$2,000:
\[
\text{Poverty Gap} = 5000 - 2000 = 3000
\]
Total Poverty Gap for these households = \( 3000 \times 20 = 60,000 \)

For the 30 households earning \$3,000:
\[
\text{Poverty Gap} = 5000 - 3000 = 2000
\]
Total Poverty Gap for these households = \( 2000 \times 30 = 60,000 \)

\end{frame}



\begin{frame}
\frametitle{Problem 3: Poverty Gap Index}


For the 50 households earning \$6,000, there is no poverty gap since income is above the poverty line.

\[
\text{Total Poverty Gap} = 60,000 + 60,000 = 120,000
\]

Total income of all households = \( 20 \times 2000 + 30 \times 3000 + 50 \times 6000 = 40,000 + 90,000 + 300,000 = 430,000 \)

Poverty Gap Index =  
\[
\text{Poverty Gap Index} = \frac{120,000}{430,000} \approx 0.279
\]
Thus, the Poverty Gap Index is approximately 0.279.

\end{frame}






\end{document}
